% Thesis chapter: conclusion and future work.
% Standalone chapter file; \input or \include from the thesis root.

\chapter{Conclusion}
\label{ch:conclusion}

\section{Summary of Findings}

This thesis set out to learn the grasp targeting decision for autonomous log
pile clearing from raw point clouds, and to carry the learned decision from
simulation onto a physical crane. Its findings form a chain.

In simulation, the exploration problem dominates the learning problem.
From-scratch reinforcement learning plateaus at roughly half the pile
regardless of observation type or reward shaping, while behavior cloning of a
privileged expert transfers full-trajectory clearing competence from raw,
unsegmented clouds, and a constrained fine-tuning recipe (frozen encoder,
privileged asymmetric critic, conservative exploration) lets on-policy RL
optimize execution without destroying that competence
(Chapter~\ref{ch:sim_study}).

Transfer is a systems problem before it is a learning problem. Centimeter
task-space tracking through the PLC interface, grapple-based camera
calibration with a floor-levelness constraint, and fused two-camera rack
localization put the real sensor data and the simulated training scene in
the same frame to within centimeters, and the gaze-phase training
environment reproduces the deployed viewpoint exactly
(Chapter~\ref{ch:sim2real}).

The action parameterization decides the failure modes. The regression policy,
as a conditional-mean estimator, averages over multimodal pile geometry and
produces an absorbing grasp-over-the-gap failure on the real machine; the
scoring policy, which classifies over observed points, is mode-seeking by
construction, keeps targets on observed material, and decouples target
selection from depth conventions. Field trials bore the analysis out: the
scoring policy cleared, without assistance, the pile shape that had forced
the heuristic baseline to be abandoned, and produced zero gap-targeting
cycles on the shape that reliably induced the regression policy's failure
(Chapter~\ref{ch:bc_policies}).

Reinforcement learning is most useful when its exploration lives in the right
space. The first-generation Gaussian fine-tuning amplified the regression
policy's mode averaging; the current architecture makes the scoring head
itself the action distribution, a categorical over observed points with
exact admissibility masking, so exploration can only propose observed
material and empty-cycle costs land on the specific point that caused them
(Chapter~\ref{ch:bcrl}).

\section{Limitations}

The task abstraction isolates targeting: motion execution is a fixed state
machine, and the simulation study removes grasped logs rather than modeling
deposition. Grasp success in simulation is approximated by post-lift
proximity rather than contact analysis. The passive-chain suspension model
affected the original stability metric strongly enough that a campaign's
stability results had to be excluded (Section~\ref{sec:sim_study_stability}),
a reminder that metric validity in simulation is itself an experimental
question; the corrected windowed metric now mirrors the on-crane
measurement. On hardware, the remaining failure classes are a fixed digging
depth that chokes the grapple on steep mounds and residual endgame
discrimination around the end of the rack; both are training-data problems
rather than architectural ones, but they are not yet solved. Finally, the
policies are deliberately not translation invariant, and deployment depends
on calibrated scene geometry with a defined recalibration procedure.

\section{Future Work}

The immediate item is completing the categorical fine-tuning campaign of
Chapter~\ref{ch:bcrl}: its evaluation battery (seeded clearing runs, the
observation-degradation sweep, offline replay on recorded trial clouds, and,
for a deployable candidate, hardware trials) is specified, with the expected
gains stated in advance on the endgame empty-cycle rate and on jagged pile
geometries.

Beyond it, three directions follow directly from the field evidence. First,
fine-tuning on real data: zero-shot transfer leaves a systematic grasp-height
bias traceable to the sim-to-real depth gap, and supervised fine-tuning on
small amounts of relabeled real-trial data is the natural correction, with
the data-efficiency curve itself a question worth answering. Second, closing
the remaining hardware failure classes: an adaptive digging depth to
eliminate chokes on steep mounds, and endgame-focused training data for the
end-of-rack region. Third, extending the learned decision beyond targeting,
toward deposition planning (structured placement into bins or trailers,
which the deployed system currently sequences with the engineered logic of
Section~\ref{sec:s2r_deposit}) and
toward richer observation pipelines; on the architecture side, the evidence
so far favors keeping the PointNet backbone and investing in the input
pipeline rather than the encoder.

The broader conclusion is methodological. On a task where exploration is
hopeless and demonstrations are cheap only in simulation, the productive
division of labor is: imitation for perception and coverage, reinforcement
learning for outcome optimization in an action space whose failure modes are
structurally bounded, and disciplined calibration so that the simulated and
physical machines disagree about as little as possible. Each stage of that
pipeline was validated on a full-scale crane, and the failures that remain
are visible, measurable, and addressable within it.
